\chapter{\(\Lambda\)-Ring Axiomatic Precision}
\label{v4-arith:gelfand-borger}

Borger's \(\Lambda\)-rings encode \(\mathbb F_1\)-descent by commuting
Frobenius lifts. The issue is whether the integral lattice underlying
\(\mathcal V^{\mathrm{prim}}\) satisfies the Wilkerson congruences.

\section{The Wilkerson Congruence and the Scalar Ring}
\label{v4-arith:gb:wilkerson}

\begin{theorem}[Wilkerson criterion, Borger 0906.3146 Prop.~1.2]
\label{v4-arith:thm:wilkerson}
\ClaimStatusProvedElsewhere. A \(\Lambda\)-ring is a commutative
ring \(R\) with commuting ring homomorphisms
\(\psi_{n}\colon R\to R\), \(n\in\mathbb{N}^{\times}\), satisfying
\(\psi_{1}=\mathrm{id}\), \(\psi_{m}\psi_{n}=\psi_{mn}\), and
\(\psi_{p}(x)\equiv x^{p}\pmod{pR}\) for every rational prime \(p\).
For torsion-free \(R\) the Wilkerson congruence is equivalent
(Wilkerson~1982) to the classical Grothendieck \(\lambda\)-ring
structure with \(\lambda^{n}\) operations.
\end{theorem}

\begin{observation}[vacuous congruence over \(\mathbb{Q}\)]
\label{v4-arith:obs:vacuous-over-Q}
Every prime \(p\) is a unit in \(\mathbb{C}\) or \(\mathbb{Q}\).
Therefore \(pR=R\) for every prime, and the Wilkerson congruence
\(\psi_{p}(x)\equiv x^{p}\pmod{pR}\) is \emph{vacuous} over
\(\mathbb{Q}\)-algebras: any two elements are congruent mod \(R\).
Under Borger's axioms, a \(\mathbb{Q}\)-algebra is automatically a
\(\Lambda\)-ring iff it carries any compatible system of commuting
ring endomorphisms \(\psi_{n}\) with \(\psi_{1}=\mathrm{id}\) and
\(\psi_{m}\psi_{n}=\psi_{mn}\).
\end{observation}

\begin{remark}
This weaker structure is what Grothendieck (Rev.~Matem.~Hispanoamer.~1971)
called a \emph{\(\psi\)-ring} or \emph{Adams-ring}. The claim
``\(\mathcal{V}^{\mathrm{prim}}\) is a \(\Lambda\)-ring'' is either
(a) trivially true over \(\mathbb{Q}\) and mathematically empty, or
(b) a claim about an \(\mathbb{F}_{1}\)- or \(\mathbb{Z}\)-integral form
\(\mathcal{V}^{\mathrm{prim}}_{\mathbb{Z}}\subset\mathcal{V}^{\mathrm{prim}}\)
on which \(p\) is not invertible, which must be specified.
\end{remark}

\section{Explicit Computations: \(\psi_{2}\) on Low-Weight States}
\label{v4-arith:gb:explicit-psi-2}

The Wilkerson congruence requires a candidate Adams operation \(\psi_{n}\)
on the Heisenberg Fock shadow of \(\mathcal{V}^{\mathrm{prim}}\). The full
\(\mathcal{V}^{\mathrm{prim}}\) is a Hochschild trace class; the Fock
calculation below takes place in the arithmetic Heisenberg sector. For
each prime \(p\), the operators
\(a^{(p)}_{-k}\) (\(k\geq 1\)) create a rank-\(1\) harmonic oscillator mode
of frequency \(k\), acting on the vacuum \(|0\rangle\). The algebra is
\([a^{(p)}_m, a^{(p)}_n]=m\delta_{m+n,0}\). On single-mode states, the
candidate Adams operation is
\begin{equation}
\label{v4-arith:eq:Tn-formula}
T_{n}\bigl(a^{(p)}_{-k}\bigr)=\sum_{d\mid n}d\cdot a^{(p)}_{-nk/d}\mathbf{1}_{d\mid k},
\end{equation}
where \(\mathbf{1}_{d\mid k}=1\) if \(d\) divides \(k\) and \(0\) otherwise;
the indicator ensures the mode index \(nk/d\) is an integer. This formula is
extended to products by the ring-homomorphism axiom.

\begin{computation}
\label{v4-arith:comp:psi2-a2}
\(\psi_{2}(a^{(2)}_{-1}|0\rangle)=a^{(2)}_{-2}|0\rangle\).
\end{computation}

\begin{computation}
\label{v4-arith:comp:psi2-a3}
\(\psi_{2}(a^{(3)}_{-1}|0\rangle)=a^{(3)}_{-2}|0\rangle\).
\end{computation}

\begin{computation}
\label{v4-arith:comp:psi2-a2a3}
Under the ring-homomorphism axiom:
\begin{equation}
\psi_{2}(a^{(2)}_{-1}a^{(3)}_{-1}|0\rangle)
\;=\;a^{(2)}_{-2}a^{(3)}_{-2}|0\rangle.
\end{equation}
\end{computation}

\begin{lemma}[composition law]
\label{v4-arith:lem:composition-law}
\ClaimStatusProvedHere. For \(m,n\geq 1\) and on single-generator
states, \(\psi_{m}\psi_{n}=\psi_{mn}\).
\end{lemma}

\begin{proof}
Compute \(\psi_{3}\psi_{2}(a^{(2)}_{-1}|0\rangle)=\psi_{3}(a^{(2)}_{-2}|0\rangle)=a^{(2)}_{-6}|0\rangle\);
and \(\psi_{6}(a^{(2)}_{-1}|0\rangle)=a^{(2)}_{-6}|0\rangle\).
Agreement is Macdonald's identity
\(p_{m(nk)}=p_{(mn)k}\) (Macdonald~I~\S8).
\end{proof}

\section{The Wilkerson Congruence Fails on the Fock Lattice}
\label{v4-arith:gb:wilkerson-fails}

\begin{computation}[Wilkerson test]
\label{v4-arith:comp:wilkerson-test}
Test \(\psi_{2}(x)\equiv x^{2}\pmod{2}\) at \(x=a^{(2)}_{-1}|0\rangle\).
\begin{itemize}
\item Left side: \(\psi_{2}(a^{(2)}_{-1}|0\rangle)=a^{(2)}_{-2}|0\rangle\).
\item Right side: \((a^{(2)}_{-1}|0\rangle)^{2}=a^{(2)}_{-1}a^{(2)}_{-1}|0\rangle\)
(normally-ordered product).
\item Difference:
\(\psi_{2}(x)-x^{2}=a^{(2)}_{-2}|0\rangle-a^{(2)}_{-1}a^{(2)}_{-1}|0\rangle
=-2L^{(2)}_{-2}|0\rangle+a^{(2)}_{-2}|0\rangle\) (using
\(L_{-2}|0\rangle=\tfrac{1}{2}a_{-1}a_{-1}|0\rangle\) for rank-\(1\)
Heisenberg). This is a specific nonzero vector.
\end{itemize}
For the congruence to hold, this vector must lie in
\(2\cdot\mathcal{V}^{\mathrm{prim}}_{\mathbb{Z}}\) for some
\(\mathbb{Z}\)-integral form.
\end{computation}

\begin{observation}[lattice-form failure]
\label{v4-arith:obs:lattice-form-failure}
On the naive Fock lattice
\(\mathcal{H}^{(p)}_{\mathbb{Z}}=\mathbb{Z}\langle a^{(p)}_{-n_{1}}\cdots a^{(p)}_{-n_{k}}|0\rangle_{p}\rangle\),
the vector \(a^{(p)}_{-2}|0\rangle-a^{(p)}_{-1}a^{(p)}_{-1}|0\rangle\)
is a primitive integer-coefficient vector; its reduction mod~\(2\)
is nonzero. The Wilkerson congruence \emph{fails} on the naive
lattice integral form.
\end{observation}

\section{The Macdonald Lattice vs the Fock Lattice}
\label{v4-arith:gb:macdonald-vs-fock-lattice}

\begin{theorem}[Macdonald I.8, Example~24]
\label{v4-arith:thm:macdonald-lattice}
\ClaimStatusProvedElsewhere. The ring \(\Lambda_{\mathbb{Z}}=
\mathbb{Z}[e_{1},e_{2},\ldots]=\mathbb{Z}[h_{1},h_{2},\ldots]\) of
symmetric functions is a \(\Lambda\)-ring satisfying the Wilkerson
congruence: the Newton identity
\(\psi_{p}(h_{n})\equiv h_{n}^{p}\pmod{p}\) in
\(\Lambda_{\mathbb{Z}}\).
\end{theorem}

\begin{observation}[Fock-Macdonald isomorphism over \(\mathbb{Q}\)]
\label{v4-arith:obs:fock-macdonald-Q}
The map \(\iota\colon\Lambda_{\mathbb{Q}}\xrightarrow{\sim}\mathrm{Fock}_{\mathbb{Q}}\),
\(p_{n}\mapsto a_{-n}|0\rangle\), is an isomorphism of
\(\mathbb{Q}\)-algebras: under the Heisenberg normalisation
\([a_m,a_n]=m\delta_{m+n,0}\) the Fock inner product
\(\langle a_{-n}|0\rangle,a_{-n}|0\rangle\rangle=n\) matches the Hall
inner product \(\langle p_n,p_n\rangle=n\), so \(\iota\) is an
isometry. Adams operations on \(\Lambda_{\mathbb{Q}}\) transport to
Fock; the \(\mathbb{Q}\)-algebra structure is preserved.
\end{observation}

\begin{observation}[lattice mismatch over \(\mathbb{Z}\)]
\label{v4-arith:obs:lattice-mismatch-Z}
Over \(\mathbb{Z}\), the two integral forms \(\iota(\Lambda_{\mathbb{Z}})\)
and the Fock lattice
\(\mathrm{Fock}_{\mathbb{Z}}=\mathbb{Z}\langle a_{-\vec{n}}|0\rangle\rangle\)
are \emph{different}. The Schur function
\(s_{2}=\tfrac{1}{2}(p_{1}^{2}+p_{2})\in\Lambda_{\mathbb{Z}}\) maps to
\(\iota(s_2)=\tfrac{1}{2}(a_{-1}^{2}+a_{-2})|0\rangle\), a half-integer
state outside \(\mathrm{Fock}_{\mathbb{Z}}\). Conversely, every Fock
monomial \(a_{-\lambda}|0\rangle=\iota(p_\lambda)\) lies in
\(\iota(\Lambda_{\mathbb{Z}})\) (power sums are integer combinations of
Schur functions, with character coefficients
\(\chi^\lambda_\mu\in\mathbb{Z}\)); the Schur basis
\(\{\iota(s_\lambda)\}\) of \(\iota(\Lambda_{\mathbb{Z}})\) generates a
strictly larger \(\mathbb{Z}\)-lattice in \(\mathrm{Fock}_{\mathbb{Q}}\)
than the Fock monomial basis, carrying denominators dividing the order
of the centralizer \(z_\mu\) at each conjugacy type \(\mu\).
\end{observation}

\section{Distinction: \(\Lambda\)-Ring vs \(\Lambda\)-Module}
\label{v4-arith:gb:ring-vs-module}

\begin{definition}[\(\Lambda\)-module, Borger 0801.1691 \S1.2]
\label{v4-arith:def:lambda-module}
A \(\Lambda\)-structure on an abelian group \(M\) is an action of the
Adams monoid \(\mathbb{N}^{\times}\) by group endomorphisms
\(\psi_{n}\). No ring structure, no congruence.
\end{definition}

\begin{theorem}[corrected form of \(\mathcal{V}^{\mathrm{prim}}\)]
\label{v4-arith:thm:corrected-lambda-form}
\ClaimStatusProvedHere. The Heisenberg Fock shadow of
\(\mathcal{V}^{\mathrm{prim}}\) over \(\mathbb{Q}\) carries a
\(\psi\)-ring structure with vacuous Wilkerson congruence via the
Hecke--Adams operations. The full
\(\mathcal{V}^{\mathrm{prim}}\) carries only the corresponding
\(\mathbb{N}^{\times}\)-module datum until a multiplication is specified.
On the natural Fock integral form there is no Borger
\(\Lambda\)-ring structure.
\end{theorem}

\section{The Log-\(p\) Rescaling Destroys Integrality}
\label{v4-arith:gb:log-p-destroys-integrality}

\begin{observation}
\label{v4-arith:obs:logp-destroys-integrality}
When the Hamiltonian is rescaled so that \(L_{0}\) assigns eigenvalue
\(k\log p\) to each weight-\(k\) state \(a^{(p)}_{-k}|0\rangle\),
the resulting grading takes values in \(\mathbb{R}\setminus\mathbb{Q}\)
(since \(\log p\) is transcendental over \(\mathbb{Q}\) for every prime
\(p\)). A \(\mathbb{Z}\)-lattice in the Heisenberg Fock shadow, spanned
by eigenstates of the conformal \(L_{0}\), is not preserved
under this rescaled grading.
\end{observation}

\begin{convention}[two-grading principle]
\label{v4-arith:conv:two-grading}
Resolve the conflict by carrying \emph{two} gradings on the same
Heisenberg basis:
\begin{enumerate}
\item \textbf{Conformal grading:} integer weight \(n\) on
\(a^{(p)}_{-n}\). Character \(\prod_{p}\prod_{n}(1-q^{n})^{-1}\)
(Heisenberg direct product, not an \(L\)-function).
\item \textbf{HY-zeta grading:} weight \(n\log p\) on
\(a^{(p)}_{-n}\). Character
\(\prod_{p}\prod_{n}(1-p^{-ns})^{-1}=\prod_{n}\zeta(ns)\) as a
formal Dirichlet series.
\end{enumerate}
The algebraic / \(\Lambda\)-module structure lives on (i); the
zeta-character lives on (ii). These are two different gradings on
the same underlying basis.
\end{convention}

\begin{remark}
This matches Bost--Connes / Connes--Consani: the \(\log p\) enters
as spectral parameter in the Hamiltonian, not as an algebraic
scalar in the underlying ring. The Bost--Connes KMS partition
function \(\zeta(\beta)\) is computed at grading (ii) on a purely
algebraic structure that lives at grading (i).
\end{remark}

\section{Is there a \(\Lambda\)-VA?}
\label{v4-arith:gb:is-there-lambda-VA}

\begin{question}
Can we define a category of \(\Lambda\)-vertex-algebras: VAs with
commuting ring-endomorphism \(\psi_{n}\) compatible with the OPE
and the vacuum?
\end{question}

\begin{proposition}[\(T_{n}\) is not a VA endomorphism]
\label{v4-arith:prop:T-n-not-VA-endo}
\ClaimStatusProvedHere. \(T_{n}\) scales conformal weight by \(n\),
i.e., \(T_{n}(a^{(p)}_{-1}|0\rangle)\) has weight~\(n\) while
\(a^{(p)}_{-1}|0\rangle\) has weight~\(1\). A VA endomorphism must
preserve conformal weight (commute with \(L_{0}\)). Therefore
\(T_{n}\) is not a VA endomorphism for \(n>1\).
\end{proposition}

\begin{theorem}[corrected category: finite Hecke-equivariant VA shadows]
\label{v4-arith:thm:HVA-category}
\ClaimStatusProvedHere. Define
\begin{equation*}
\mathrm{HVA}=\{(V,\{T_{n}\}_{n\in\mathbb{N}^{\times}})\colon
V\text{ VA}, T_{n}\text{ linear, }T_{m}T_{n}=T_{mn},
T_{n}|0\rangle=|0\rangle\}.
\end{equation*}
For every finite set \(S\) of primes, the finite Heisenberg truncation
\(\mathcal{H}^{(S)}=\bigotimes_{p\in S}\mathcal{H}^{(p)}\) is an object
of \(\mathrm{HVA}\), with the further intertwining property
\(T_n(a^{(p)}_{-1}|0\rangle)=a^{(p)}_{-n}|0\rangle\) on level-1
generators (\Cref{v4-arith:prop:T-n-not-VA-endo}). On higher-weight
states \(T_n\) produces sums of mixed weight, e.g.\
\(T_2(a^{(p)}_{-2}|0\rangle)=a^{(p)}_{-4}|0\rangle+2\,a^{(p)}_{-2}|0\rangle\),
so neither \([T_n,L_{0}]\) nor \([T_n,L_{-1}]\) vanishes as an operator
identity; direct computation gives
\([T_2,L_{-1}]\,a^{(p)}_{-1}|0\rangle=a^{(p)}_{-4}|0\rangle+2\,a^{(p)}_{-2}|0\rangle-2\,a^{(p)}_{-3}|0\rangle\neq 0\).
The category \(\mathrm{HVA}\) records the Adams multiplicativity
\(T_mT_n=T_{mn}\) and the vacuum preservation; the conformal
interaction is a separate datum. The restricted product is an
\(\mathbb{N}^{\times}\)-module shadow, not a vertex algebra structure on
\(\mathcal{V}^{\mathrm{prim}}\).
\end{theorem}

\begin{corollary}[\(\psi\)-ring plus Hecke-equivariance]
\label{v4-arith:cor:psi-ring-hecke}
\ClaimStatusProvedHere. The correct categorical statement:
\(\mathcal{V}^{\mathrm{prim}}\in\mathrm{Mod}_{\mathbb{N}^{\times}}(\mathrm{Vect}_{\mathbb{C}})\)
with the \(\mathbb{N}^{\times}\)-action by Hecke operators \(T_{n}\),
whose finite Heisenberg shadows are the Hecke-equivariant VA objects of
\Cref{v4-arith:thm:HVA-category}. This is Mackey's Hecke algebra
representation enhanced by finite Fock shadows, not Borger's
\(\Lambda\)-ring.
\end{corollary}

\section{The Descent Interpretation Goes Through Bost--Connes}
\label{v4-arith:gb:descent-bost-connes}

\begin{remark}
Borger's \(\mathbb{F}_{1}\)-descent programme (arXiv:0906.3146~\S3)
characterizes a \(\Lambda\)-structure on \(R\) as a descent datum from \(R\) to
\(R\otimes_{\mathbb{F}_{1}}\mathbb{Z}\). The \(\mathbb{F}_{1}\)-descent of
\(\mathcal{V}^{\mathrm{prim}}\) in Borger's sense requires a genuine
\(\Lambda\)-ring structure, hence a Wilkerson-congruent integral form, which
the natural Fock lattice does not supply (Theorem~\ref{v4-arith:thm:corrected-lambda-form}
and Observation~\ref{v4-arith:obs:lattice-form-failure}). The correct
arithmetic framework for \(\mathcal{V}^{\mathrm{prim}}\) is
Bost--Connes (Selecta~Math.~1 (1995), 411--457) and Connes--Consani
(arXiv:1405.4527, the arithmetic site).
\end{remark}

\begin{theorem}[Bost--Connes as the correct descent]
\label{v4-arith:thm:bost-connes-descent}
\ClaimStatusConjectured. \(\mathcal{V}^{\mathrm{prim}}\) has a
Heisenberg-Fock shadow compatible with the Bost--Connes Hecke
\(C^{*}\)-dynamical system, with:
\begin{itemize}
\item Hecke action \(\{T_{n}\}\) = Bost--Connes semigroup action;
\item \(\mathbb{Q}/\mathbb{Z}\)-Galois symmetry = Bost--Connes
class-field-theory symmetry (Bost--Connes Thm~5);
\item graded character = HY \(L\)-function after rescaling by
\(\log p\).
\end{itemize}
Borger's \(\mathbb{F}_{1}\)-descent requires a Wilkerson-congruent
integral \(\Lambda\)-form on the natural Fock lattice; the Macdonald
lattice \(\Lambda_{\mathbb{Z}}\) carries such a form but differs from
the Fock lattice over \(\mathbb{Z}\)
(\Cref{v4-arith:obs:lattice-mismatch-Z}). The Bost--Connes /
Connes--Consani arithmetic site supplies the correct framework via
tropical and semiring structures.
\end{theorem}

\begin{remark}[the three-level structure]
\label{v4-arith:rem:three-level-structure}
The Wilkerson congruence is vacuous over \(\mathbb{Q}\)
(\Cref{v4-arith:obs:vacuous-over-Q}), so the distinction between
\(\Lambda\)-ring and \(\psi\)-ring is a purely integral question.
The Macdonald lattice \(\Lambda_{\mathbb{Z}}\) carries the Wilkerson
congruence (\Cref{v4-arith:thm:macdonald-lattice}); the Fock lattice
\(\mathbb{Z}\langle a^{(p)}_{-\vec{n}}|0\rangle\rangle\) does
not (\Cref{v4-arith:obs:lattice-form-failure}), and the two integral
forms are distinct (\Cref{v4-arith:obs:lattice-mismatch-Z}). At finite
Heisenberg level, \(T_{n}\) scales conformal weight
(\Cref{v4-arith:prop:T-n-not-VA-endo}), so the category \(\mathrm{HVA}\)
applies only to finite Fock shadows
(\Cref{v4-arith:thm:HVA-category}). The \(\mathbb{F}_{1}\)-descent
framework adequate to \(\mathcal{V}^{\mathrm{prim}}\) is
Bost--Connes rather than Borger (\Cref{v4-arith:thm:bost-connes-descent}).
\end{remark}
